% 第 12 周结项演示视频讲稿（口语版·PPT＋实验结果＋报告＋代码多窗口）——由幻灯同一份源生成，数字取自 docs/week12/numbers.tex；编译：XeLaTeX
\documentclass[11pt]{ctexart}
\usepackage[a4paper, margin=2.3cm]{geometry}
\usepackage{amssymb}
\usepackage{url}
\usepackage{enumitem}
\setlist{nosep}
\usepackage{parskip}
\usepackage{xcolor}
\newcommand{\cue}[1]{\noindent\colorbox{gray!15}{\parbox{\dimexpr\linewidth-2\fboxsep}{\textbf{#1}}}\par\smallskip}

\title{\textbf{结项演示视频完整讲稿（口语版）}\\
\large Chooser Option 定价项目 · 第 1--12 周 · 10--15 分钟 · 对着 PPT＋实验结果＋报告＋代码讲}
\author{佟铸城（Zhucheng Tong）}
\date{2026 年 10 月}

\begin{document}
\maketitle

\section*{录制前准备（约 10 分钟）}
形式与第 8 周相同：在 PPT、实验结果、报告和代码之间切换。提前打开并排好下列窗口（用多桌面或 Cmd+Tab 切换）：
\begin{enumerate}
  \item \textbf{PPT}：\texttt{docs/week12/final\_presentation.pptx}，用 PowerPoint 或 Keynote \textbf{全屏放映}（15 页；第 8 周录的是编辑视图，这次建议放映模式）；
  \item \textbf{工具}：先开 Docker，在项目目录运行 \texttt{docker compose up -d}，浏览器打开 \texttt{localhost:8501}，放大窗口；左栏“标的”默认就是 JPM，先把“定价”“全市场概览”“运行状态”三个页签各点一遍，让数据加载好；
  \item \textbf{报告}：\texttt{docs/week12/final\_report\_cn.pdf}，提前翻到第 13 页（逐年柱状图）；第 14 页是各锚误差图（可选）；
  \item \textbf{代码}：编辑器打开 \path{src/models/chooser.py}、\path{src/ml/cross_sectional.py} 和 \path{docs/week12/preregistration.json}；
  \item \textbf{录屏}：按 Cmd+Shift+5，选“录制整个屏幕”，在“选项”里勾上麦克风；开专注模式屏蔽通知；接上电源，避免中途睡眠。
\end{enumerate}
先不录音干跑一遍，熟悉切屏路线；可以分段录（每段一到三页），录完交给我无损拼接。

\section*{讲稿（总时长约 13--14 分钟）}

\cue{[0:00--0:50 · PPT 第 1--2 页（标题＋项目一览）]}
老师好，我是佟铸城。这个视频是 chooser option 定价项目的结项汇报，覆盖全部十二周；前八周用三页快速带过，重点讲后四周。先看四个数字。前两个来自前八周：论文定价复现出来是 29.13 美元；BSM“波动率不变”这个假设的代价，大约是价格的 35\%。后两个来自后四周：在 2023 到 2024 年的样本外、505 只股票上，把波动率的“锚”从 EWMA 换成 GARCH，每份合约平均少错 1.07 美元；机器学习在 GARCH 锚之上再少错 0.337 美元，95\% 区间 0.206 到 0.456，检验判据在跑数之前就登记好了。第八周的结论是机器学习在 JPM 一只股票上没跑赢基准；扩到 S\&P 500 之后，第一次在样本外得到了有边界的正增量。

\cue{[0:50--2:00 · PPT 第 3 页 → 切到代码 src/models/chooser.py（拆分函数＝第 19--26 行；教材公式版＝第 29--42 行，一瞥即回）]}
先回顾定价本身。Chooser option 买的时候不用定方向，到选择日再决定是看涨还是看跌。用 put-call 平价关系变形以后，它可以拆成一个到期日在 T2 的 call，加一个到期日在选择日、行权价折算过的 put，所以能直接写出定价公式。（切到 chooser.py）代码里就是这个不到十行的拆分函数；下面是照教材公式独立写的另一版，两份互相对答案，一致到小数点后十五位。（回 PPT）价格大小关系也对：比只买一个期权贵，比同时买 call 和 put 便宜。

\cue{[2:00--3:10 · PPT 第 4 页]}
基准测试的思路是事后对答案：同一份合约，用当时能看到的波动率定一次价，再用后来一年真实发生的波动率定一次，差距就是“波动率不变”的代价——最好的历史估法也有 35\%。机器学习只学“相对基准该往上还是往下调多少”，但测试集打开以后，没有一个模型赢过基准：基准误差 8.51 美元，最好的机器学习 14.75 美元。主要原因是一只股票的历史里，可学的完整市场周期太少；后四周导师给的方向之一，正好是把数据扩到更大的股票池。

\cue{[3:10--3:50 · PPT 第 5 页]}
后四周，导师给了三个工程方向：让数据管道和工具在线运行；把标的扩到更大的股票池；测量全流程的时间和空间占用并优化。另外还有一条理论线：在几百只股票上重新检验“锚”和机器学习，最后回答您的两个问题——机器学习增量的显著性边界，以及隐含波动率怎么定位。

\cue{[3:50--5:00 · PPT 第 6 页 → 切到浏览器工具（localhost:8501）“运行状态”页签，指运行日志表 → 回 PPT]}
第九周把整套东西搬上了在线栈：数据存在 PostgreSQL，报价放在 Redis 缓存里，调度器每个交易日收盘后自动跑一遍，五个容器一条命令启动。验收标准只有一条：同一份输入，新旧两套算术给出相同结果。（切到工具“运行状态”页签）这是每一次运行的日志，跑了多少只股票、每个阶段多少秒，都记在库里。第十二周修了一个运维问题：笔记本睡眠时虚拟机被挂起，调度器的计时跟着停，9 月 25 号之后就没再跑过；现在改成每 10 分钟一次心跳，漏跑就立即补跑。（回 PPT）

\cue{[5:00--5:50 · PPT 第 7 页]}
第十周扩到 S\&P 500。关键是“点时成分”：用成分变更记录倒推出每只股票哪天进、哪天出指数，面板每一天只用当天的成员，不拿今天的成分股回看过去。面板一共 617,211 行、532 只股票。要说明的局限是：有 51 只成员退市太早、数据源不再提供历史，所以面板有生存偏差。面板上的第一项发现：GARCH 1 日预测作锚，在 72\% 的股票上胜过 EWMA，成了下一轮的候选锚。

\cue{[5:50--6:50 · PPT 第 8 页]}
第十一周做性能，先剖析、再优化，同一个脚本、同一台机器前后对比。剖析发现报价慢在读数据、不在计算，日更有一半时间在排队等网络请求。结果：全池报价从 390 毫秒降到 23.5 毫秒，标的数还从 578 只变成了 1496 只；市场数据阶段从 12.8 秒降到 2.6 秒。浅蓝这一行是被测量否定的优化：往返次数大降，时间几乎没变，瓶颈不在这里——代码留着，结论如实记录。压测到 S\&P 1500，1517 只的稳态日更是 59.5 秒。

\cue{[6:50--8:00 · PPT 第 9 页 → 切到代码 src/ml/cross\_sectional.py 文件头（第 4--7 行、第 15--16 行），回 PPT]}
下面是理论线。第一轮面板机器学习在验证集上全面失败，原因是训练期和验证期之间，整个面板“该往上还是往下调”的幅度一起漂移了。所以第二轮只学“横截面”：每一天内部，哪些股票该比锚调高、哪些该调低；锚也换成了 GARCH。测试集只打开一次：以 EWMA 为 1，单用 GARCH 锚是 0.791，加上机器学习是 0.755——大部分改进来自换锚，机器学习再进一步。拆开检查：只保留同一天内的重新排序是 0.742，只保留整体水平的平移是 0.801，比不修正还差——收益全部来自横截面。（切到 cross\_sectional.py）部署的就是这个模型：目标按日期去均值，特征在每个日期内标准化，定价时用锚乘上预测的修正倍数。（回 PPT）

\cue{[8:00--9:10 · PPT 第 10 页 → 切到 docs/week12/preregistration.json（滚动一屏即回）]}
您的第一个问题是机器学习增量的显著性边界。我最在意的是不让自己挑结果，所以先把判据写进这个文件，（切到 preregistration.json）记下 SHA-256 指纹再跑数，审计会核对指纹。（回 PPT）第二步是检验“检验程序本身”：同行业的股票会一起涨跌，按单只股票重抽样会让区间偏窄。所以每个假设先在“真实增量为零”的模拟数据上试，名义 5\% 的检验，实际拒绝率要在 2.5\% 到 7.5\% 之间才用。按股票抽样在四个假设上全部不合格，改为先抽行业、再抽股票；这个检验工具本身也在已知答案的数据上验证过，强行业冲击下按股票抽样的拒绝率升到 69\%，被正确识别出来。最后四个假设一起做 Holm 矫正。

\cue{[9:10--10:20 · PPT 第 11 页 → 切到报告 docs/week12/final\_report\_cn.pdf 第 13 页（逐年柱状图），回 PPT]}
结果在这张表里。主假设 H2，机器学习增量每份合约 0.337 美元，95\% 区间 0.206 到 0.456，按事先规则判定成立。但边界要说清楚。（切到报告第 13 页）逐年拆开看，增量主要来自 2024 年，0.634 美元，2023 年只有 0.048；独立年份只有两个，所以只能说这段时期内它普遍存在于各只股票，不能外推到别的年份。换锚 H1 的增量是 1.07 美元，约为机器学习的三倍、每年都为正，但两种检验程序都不合格，只作描述——这不推翻换锚，因为训练、验证、测试三个窗口它都占优。（回 PPT）

\cue{[10:20--11:30 · PPT 第 12 页（可选：切到报告第 14 页的各锚误差图）]}
第二个问题是隐含波动率。数据是学校 WRDS 的 OptionMetrics，按许可只在本机研究。在 506 只股票的共同样本上，它是最强的锚：平均误差 5.16 美元，在 92\% 的股票上胜过 GARCH 1 日锚。但它强在水平、不在择时：和事后真值的横截面相关是 0.95，时间上的相关中位数只有 0.11。按事先的判据，它定位为研究基准、不进工具：对 GARCH 1 年锚的优势，检验程序不合格；而且 OptionMetrics 只到 2025 年 8 月底，公开期权链只有当天快照，可部署性无法核实。顺带更正第十一周：当时 JPM 上的“相关 0.78”是机械相关，把隐含波动率换成常数反而是 0.83。

\cue{[11:30--12:40 · PPT 第 13 页 → 切到浏览器工具“定价”页签（左栏标的为 JPM）（指主定价标签里的“GARCH 1 日”、ML 参考价及下方提示）→“全市场概览”（指“锚”列）→ 回 PPT]}
这些结论落到了工具上。（切到工具“定价”页签）主定价换成了 GARCH 1 日锚：每只股票每天存参数和状态，日更接着往下算，只在排定的日子重估。历史不足 750 个交易日的新股退回 EWMA，每笔报价都标明用的哪个锚——10 月 2 日的 1,493 笔里，15 笔用 EWMA。误差区间是平静 5.48、恐慌 5.86 美元，但校准窗口不含 2020 年崩盘，危机中可能偏窄，界面已注明。机器学习参考价换成了刚才检验的模型，覆盖 496 只 S\&P 500 成员，仍只作参考。（切到“全市场概览”）这一列能看到每只股票用的锚。EWMA 主定价一个环境变量就能切回；之前各周的审计全部保留并通过，第十二周新增 12 项。（回 PPT）

\cue{[12:40--13:40 · PPT 第 14--15 页（结论＋谢谢）]}
最后回到您最早提的“1+1”的问题。十二周下来，证据指向同一个结构：随机过程和机器学习是分工，不是比赛。同一组样本上，换锚的改进 1.07 美元，约为机器学习 0.337 美元的三倍——GARCH 管波动率的水平，机器学习管同一天里股票之间的相对差异。隐含波动率的价值同样在水平定位，眼下只作研究基准。接下来值得做三件事：用隐含波动率的水平设定 GARCH 的长期方差；积累更长的样本外历史；美式 chooser 和机器学习定价加速器。（翻到最后一页）所有结果固定了随机种子，可以从头复现，代码、逐周报告和审计脚本都在仓库里。谢谢老师。

\section*{录后检查清单}
\begin{itemize}
  \item[$\square$] 总时长 10--15 分钟（超时可压缩第 3--4 页回顾段或第 8 页性能段）
  \item[$\square$] 回放开头 30 秒确认有声音；无通知弹窗、无个人信息入镜（注意终端与浏览器标签页）
  \item[$\square$] 工具画面里没有出现 WRDS 相关的任何文件或窗口
  \item[$\square$] 文件命名 \texttt{chooser\_demo\_w12\_佟铸城.mov}；分段录制或文件过大再找我拼接、压缩
\end{itemize}
\end{document}
